\documentclass[a4paper,12pt]{article}

%----------------------------------------------------------------------------------------
%	PACKAGES
%----------------------------------------------------------------------------------------
\usepackage{url}
\usepackage{parskip} 	

%other packages for formatting
\RequirePackage{color}
\RequirePackage{graphicx}
\usepackage[usenames,dvipsnames]{xcolor}
\usepackage[scale=0.9]{geometry}

%tabularx environment
\usepackage{tabularx}

%for lists within experience section
\usepackage{enumitem}

% centered version of 'X' col. type
\newcolumntype{C}{>{\centering\arraybackslash}X} 

%to prevent spillover of tabular into next pages
\usepackage{supertabular}
\usepackage{tabularx}
\newlength{\fullcollw}
\setlength{\fullcollw}{0.47\textwidth}

%custom \section
\usepackage{titlesec}				
\usepackage{multicol}
\usepackage{multirow}

\titleformat{\section}{\Large\scshape\raggedright}{}{0em}{}[\titlerule]
\titlespacing{\section}{0pt}{10pt}{10pt}

\usepackage[style=authoryear,sorting=ynt, maxbibnames=2]{biblatex}

\usepackage[unicode, draft=false]{hyperref}
\definecolor{linkcolour}{rgb}{0,0.2,0.6}
\hypersetup{colorlinks,breaklinks,urlcolor=linkcolour,linkcolor=linkcolour}
\addbibresource{citations.bib}
\setlength\bibitemsep{1em}

\usepackage{fontawesome5}


% job listing environments
\newenvironment{jobshort}[2]
    {
    \begin{tabularx}{\linewidth}{@{}l X r@{}}
    \textbf{#1} & \hfill &  #2 \\[3.75pt]
    \end{tabularx}
    }
    {
    }

\newenvironment{joblong}[2]
    {
    \begin{tabularx}{\linewidth}{@{}l X r@{}}
    \textbf{#1} & \hfill &  #2 \\[3.75pt]
    \end{tabularx}
    \begin{minipage}[t]{\linewidth}
    \begin{itemize}[nosep,after=\strut, leftmargin=1em, itemsep=3pt,label=--]
    }
    {
    \end{itemize}
    \end{minipage}    
    }



%----------------------------------------------------------------------------------------
%	BEGIN DOCUMENT
%----------------------------------------------------------------------------------------
\begin{document}

% non-numbered pages
\pagestyle{empty} 

%----------------------------------------------------------------------------------------
%	TITLE
%----------------------------------------------------------------------------------------


\begin{tabularx}{\linewidth}{@{} C @{}}
\Huge{Ian Chen} \\[7.5pt]
\href{https://github.com/ianchu0317}{\raisebox{-0.05\height}\faGithub\ ianchu0317} \ $|$ \ 
\href{https://linkedin.com/in/ianchenn}{\raisebox{-0.05\height}\faLinkedin\ ianchenn} \ $|$ \ 
\href{https://ianchenn.com}{\raisebox{-0.05\height}\faGlobe \ ianchenn.com} \ $|$ \ 
\href{mailto:chenian317@gmail.com}{\raisebox{-0.05\height}\faEnvelope \ chenian317@gmail.com}  \\
\end{tabularx}

%----------------------------------------------------------------------------------------
% ABOUT
%----------------------------------------------------------------------------------------

%Interests/ Keywords/ Summary
\section{Sobre mi}
Estudiante de Ingeniería en Informática (FIUBA) con interés en backend, infraestructura, administración de sistemas y automatización. Experiencia en proyectos personales y académicos donde integré Linux, Docker y Python. Busco desarrollarme en roles técnicos como Backend o SysAdmin Junior para aplicar y seguir aprendiendo en entornos reales.

%----------------------------------------------------------------------------------------
% EXPERIENCE SECTIONS
%----------------------------------------------------------------------------------------

\section{Experience}

\begin{tabularx}{\linewidth}{ @{}l r@{} }
\textbf{Administrador de Redes y Sistemas} \\[3.75pt]
\multicolumn{2}{@{}X@{}}{
- Instalar, mantener, configurar y actualizar servidores de Linux con virtualización en Proxmox. 

- Deploy proyectos y servicios contenerizados en Docker. 

- Documentación en Bookstack y GitHub. 

- Desarrollar scripts de Bash y Crontab para Backup automatizados.  

- Tecnologías: Proxmox, Linux, Debian, Ubuntu, Docker, Tailscale, Cloudflare, Crontab, Backups, SSH.}  \\
\end{tabularx}

%----------------------------------------------------------------------------------------
% PROJECTS
%----------------------------------------------------------------------------------------

%Projects
\section{Proyectos}

\begin{tabularx}{\linewidth}{ @{}l r@{} }
\textbf{Fobium - Foro web} & \hfill \href{https://www.fobium.com}{fobium.com} \\[3.75pt]
\multicolumn{2}{@{}X@{}}{
Página web de un foro anónimo para publicar fobias. Mi rol en este trabajo grupal fue el desarrollo de Backend y organización del grupo y el proyecto en general. Tecnologías utilizadas: Python, FastAPI, Postgres SQL, Docker, {Gemini AI}, Git/GitHub, Cloudflare, Proxmox, Trello}  \\
\end{tabularx}


\begin{tabularx}{\linewidth}{ @{}l r@{} }
\textbf{ToDoList} & \hfill \href{https://todo.ianchenn.com}{todo.ianchenn.com} \\[3.75pt]
\multicolumn{2}{@{}X@{}}{
Página web para organizar lista de tareas con autenticacion JWT. Tecnologías utilizadas: FastAPI, MySQL, Docker, Git, Cloudflare, Proxmox 
}  \\
\end{tabularx}


\begin{tabularx}{\linewidth}{ @{}l r@{} }
\textbf{Led API} & \hfill \href{https://github.com/ianchu0317/LED-control-API}{GitHub} \\[3.75pt]
\multicolumn{2}{@{}X@{}}{
Patrones LED controlados por medio de API e interfaz en página web. Tecnologías utilizadas: Arduino, Raspberry Pi, FastAPI, C++, Python
}  \\
\end{tabularx}


\begin{tabularx}{\linewidth}{ @{}l r@{} }
\textbf{Blog de Matemática} & \hfill \href{https://ianchu0317.github.io/analisis-matematico-ii/}{analisis-ii.com} \\[3.75pt]
\multicolumn{2}{@{}X@{}}{
Blog de resoluciones y parciales para análisis matemático 2 UBA. Tecnologías utilizadas: Github Pages, LaTeX, Google Analytics, Google Search Console
}  \\
\end{tabularx}

%----------------------------------------------------------------------------------------
%	SKILLS
%----------------------------------------------------------------------------------------
\section{Aptitudes}
\begin{tabularx}{\linewidth}{@{}l X@{}}
Idiomas &  \normalsize{Ingles, Español, Chino Mandarín.}\\
Lenguajes  &  \normalsize{Python, Go, C, Bash, Markdown, LaTeX.}\\ 
Sistemas  \& herramientas &  \normalsize{Windows, Linux, Ubuntu, Debian, Proxmox, SSH, Systemd, Cron.}\\ 
Redes  &  \normalsize{Configuración de router, optimización de cableado, testeo de velocidad.}\\ 
Otros  &  \normalsize{Gestión de proyectos, Documentación, Code review, Code testing, Instalación de sistemas.}\\ 
\end{tabularx}

%----------------------------------------------------------------------------------------
%	EDUCATION
%----------------------------------------------------------------------------------------
\section{Educacion}
\begin{tabularx}{\linewidth}{@{}l X@{}}	
Ingeniería Informática & Facultad de Ingeniería UBA \textbf{(FIUBA)} \hfill \normalsize  \\

Secundario Completo & Colegio Nacional de Buenos Aires \textbf{(CNBA)} \hfill (Promedio: 8.7/10.0) \\ 

\end{tabularx}

%----------------------------------------------------------------------------------------
%	PUBLICATIONS
%----------------------------------------------------------------------------------------
% \section{Publications}
% \begin{refsection}[citations.bib]
% \nocite{*}
% \printbibliography[heading=none]
% \end{refsection}


% \vfill
% \center{\footnotesize Last updated: \today}

\end{document}
